%% ---------------------------------------------------------------------------------------------------------
\section{Discussion}\label{sec:discussion}

\paragraph{What the confirmatory results establish} Two pre-registered statements hold on new scenes: the conditional harm of the deployed system is below 5\% for both policies at family-wise error 0.05 (calibration), and the system succeeds more often than the policy alone for both policies (Final Test). Everything else in this paper, including the mechanism, the per-stage decisions, the bounds restricted to acting scenes and the bootstrap intervals, is descriptive or post hoc.

\paragraph{Separating rescue from harm changes what is reported} A success rate of 94.75\% against 87.5\% does not tell a user whether the intervention ever breaks a working episode. The paired four-cell count does: on the Final Test, 1 of 697 baseline successes was lost while 56 of 103 failures were rescued (pooled, descriptive). The certificate turns the calibration counterpart of this statement into a bound with a stated error probability, although most of its margin comes from scenes in which V2 does not move.

\paragraph{Abstention as a precaution, not a demonstrated cause} V2 acted in about one in five triggered scenes (41 of 195 and 39 of 191 in development). Most triggered scenes were already aligned, and in the rest V2 abstained when an estimate was missing, uncertain or contradicted. The rules-off ablation, which removes the three V2-specific rules as one bundle, acted in 53 and 50 development scenes instead of 41 and 39, gained one and three rescues, and had the same harms (1 and 0). The development data therefore do not show that these rules reduce harm, and the ablation cannot attribute any effect to a single rule. We kept the rules as a precaution before the confirmatory stages, and the certificate holds for the system with them.

\paragraph{Harm through time and incomplete motion} The observed development harm was not a collision but a combination of a move that covered less than its commanded path and a spent time budget. V2 already closes the loop on its own estimates after every step, but it does not check before acting whether the remaining offset can be removed within its step budget at the executed rate, nor whether the steps it spends leave the policy enough time. A feasibility gate based on the observed execution rate and the remaining horizon, or an estimate of the value of acting from earlier outcomes, are possible extensions; neither was implemented or tested.

\paragraph{What the memory does} Memory let V2 act when the target was hidden at the trigger. Because the ablation that removes earlier frames also lowers the confirmation count, our data support the statement that the confirmed multi-frame estimate improves on a single-frame estimate, not that memory alone does.

\paragraph{Wrong-object triggers} The trigger is an observable rule and can fire when the policy holds the wrong object, as in scene 51152. There V2 abstained for lack of a held-cube estimate, while the privileged oracle, which trusts the true cube positions, pushed the grasped target. Abstaining on missing estimates is therefore also a guard against a premise of the trigger that the robot cannot verify.
